\documentclass[11pt,a4paper]{article}
\usepackage[margin=2.5cm]{geometry}
\usepackage{times}
\usepackage{latexsym}
\usepackage{graphicx}
\usepackage{amsmath}
\usepackage{amssymb}
\usepackage{booktabs}
\usepackage{multirow}
\usepackage{xcolor}
\usepackage{colortbl}
\usepackage{array}
\usepackage{adjustbox}
\usepackage{subcaption}
\usepackage{natbib}
\usepackage{hyperref}

\definecolor{lightgreen}{rgb}{0.85,0.95,0.85}
\definecolor{lightyellow}{rgb}{1.0,1.0,0.85}
\definecolor{lightred}{rgb}{1.0,0.85,0.85}
\definecolor{lightgray}{rgb}{0.92,0.92,0.92}

\newcommand{\caclr}{CaCLR}
\newcommand{\dauroc}{$\Delta$AUROC}

\title{Cross-Lingual Correctness Probe Transfer Reveals\\Architecture-Dependent Pareto Frontiers:\\Discovery, Negative Result, and Practical Recommendation}

\author{Anonymous}

\begin{document}
\maketitle

% ============================================================
\begin{abstract}
Probe-based calibration methods exploit internal correctness directions to assess LLM confidence, but remain English-only.
We present the first systematic evaluation of cross-lingual correctness probe transfer across 20 languages, two MoE architectures (Qwen3-235B-A22B, Llama~4 Scout), and one dense architecture (Gemma~4 31B), revealing three findings.
\textbf{(1)~Discovery:} A depth-transferability Pareto frontier emerges whose shape is architecture-dependent---Qwen3 shows a uniformly negative tradeoff (per-language Spearman $\rho \approx -0.50$) while Llama~4 Scout shows language-specific behavior (global $\rho = -0.12$; 4/20 languages exhibit \emph{positive} $\rho$).
\textbf{(2)~Principled negative result:} CaCLR, a calibration-aware probe training objective combining BCE, smooth-binned ECE, and MMD consistency loss, fails to shift this frontier: a $\beta$ sweep across three orders of magnitude yields negligible effect ($d = 0.07$, $p = 0.38$). CKA analysis suggests MMD matches marginals but cannot align the low-dimensional correctness subspace that rotates with depth.
\textbf{(3)~Practical recommendation:} Simple linear projection with ${\sim}50$ parallel pairs provides +3.5pp AUROC over zero-shot transfer, surpassing CaCLR at only 20 pairs.
A methodological warning accompanies these findings: control task analysis reveals that ${\sim}27\%$ of language-layer probing combinations exhibit selectivity below 0.05, concentrated in low-resource languages and shallow layers.
\end{abstract}

% ============================================================
\section{Introduction}

Non-English languages suffer substantially worse calibration in multilingual LLMs: expected calibration error (ECE) can be 5$\times$ higher than English.
Recent work has shown that internal \emph{correctness directions}---linearly accessible from intermediate representations---are orthogonal to verbalized confidence \cite{miao2026steering}, implying that output-based calibration is structurally limited.
Probe-based methods such as CORAL \cite{miao2026coral} and correctness-confidence steering \cite{miao2026steering} exploit this finding, but operate exclusively in English.
Whether this entire research line generalizes beyond English is an open question.

We provide the first systematic test.
Probing 20 languages across two large-scale MoE architectures (Qwen3-235B-A22B, 94 layers; Llama~4 Scout, 48 layers) and one dense architecture (Gemma~4 31B, 60 layers), we discover that cross-lingual correctness probe transfer is governed by a \emph{depth-transferability Pareto frontier} whose shape is architecture-dependent (Section~\ref{sec:pareto}).
On Qwen3, all 20 languages exhibit a consistent negative tradeoff between probe depth (and thus English AUROC) and cross-lingual transfer quality ($\rho_{\text{median}} = -0.48$).
On Llama~4 Scout, the tradeoff is language-specific: Arabic, Korean, and Hindi improve with depth while Chinese and Swahili degrade, producing a near-zero global correlation ($\rho = -0.12$).

To test whether distributional matching can shift this frontier, we design CaCLR---a calibration-aware cross-lingual regularization objective combining BCE with smooth-binned ECE ($\alpha$) and MMD consistency ($\beta$) losses.
A systematic $\beta$ sweep across three orders of magnitude yields a principled negative result: the MMD component is inert (Cohen's $d = 0.07$, $p = 0.38$; Section~\ref{sec:caclr}).
CKA geometric analysis provides a suggestive explanation: MMD matches marginal distributions but cannot align the low-dimensional directional structure of correctness representations, which rotates as a function of depth (Section~\ref{sec:cka}).

The constructive outcome is a practical recommendation: simple linear projection with ${\sim}50$ parallel pairs yields +3.5pp AUROC over zero-shot transfer, surpassing CaCLR at only 20 pairs (Section~\ref{sec:linproj}).
Difficulty-controlled decomposition reveals that the transferable component of correctness signal is difficulty-shared and language-agnostic, while the degrading component is capability-specific (Section~\ref{sec:decomp}).
Cross-lingual orthogonality between correctness and verbalized confidence persists across all 20 languages (max cosine 0.11), extending prior English-only findings (Section~\ref{sec:ortho}).

A methodological warning accompanies our findings: Hewitt \& Liang control task analysis reveals that ${\sim}27\%$ of language-layer probing combinations exhibit selectivity below 0.05, concentrated in low-resource languages and shallow layers (Section~\ref{sec:control}).
We recommend a minimum selectivity threshold for interpretable cross-lingual probing.

% ============================================================
\section{Related Work}

\paragraph{Probe-Based Calibration.}
CORAL \cite{miao2026coral} trains cross-task correctness probes showing that correctness signals transfer across benchmarks but not across languages.
Miao \& Ungar (\citeyear{miao2026steering}) demonstrate correctness-confidence orthogonality in English, establishing that output-based calibration is structurally limited.
Ji et al.\ identify internal confidence directions as a ``manipulable knob'' for calibration.
We extend these English-only findings to 20 languages.

\paragraph{Cross-Lingual Probing.}
Xiao et al.\ (\citeyear{xiao2025}) transfer \emph{knowledge boundary} probes (input-level answerability) across 8 languages.
Our work targets \emph{correctness} probes (output-dependent accuracy)---a fundamentally different signal---across 20 languages, with calibration rather than probing accuracy as the downstream metric.
LACE \cite{zhou2024lace} ensembles per-language layers but requires target-language labels.
We target zero-label transfer.

\paragraph{Cross-Lingual Alignment.}
Linear projection for cross-lingual alignment is well-established in word embeddings (MUSE, \citealp{conneau2018muse}; VecMap, \citealp{artetxe2018vecmap}).
Our linear projection is conceptually similar but operates on sentence-level correctness representations; the contribution is demonstrating effectiveness for probe transfer and characterizing sample efficiency.

\paragraph{Domain Adaptation Theory.}
Ben-David et al.\ (\citeyear{bendavid2010}) provide H-divergence bounds for domain adaptation; Zhao et al.\ (\citeyear{zhao2019}) show that marginal distribution matching fails when conditional distributions $P(Y|X)$ differ across domains.
This directly predicts CaCLR's MMD failure.

\paragraph{Representation Similarity.}
CKA \cite{kornblith2019cka} and SVCCA \cite{raghu2017svcca} measure representational alignment.
We use CKA as a correlational (not mechanistic) analysis tool, reporting partial correlations controlling for depth.

\paragraph{Probing Methodology.}
Hewitt \& Liang (\citeyear{hewitt2019control}) introduce control tasks for probe validation.
We extend their framework cross-lingually and identify systematic capacity artifacts.

% ============================================================
\section{Method}

\subsection{Correctness Probes}
\label{sec:probes}

Given a multilingual LLM with $L$ layers and an English MCQA dataset with items $\{(x_i, y_i)\}$ where $y_i \in \{0, 1\}$ indicates correctness of the model's predicted answer, we train a logistic regression probe at each layer $\ell$:
\begin{equation}
p_\ell(y=1 \mid \mathbf{h}_\ell) = \sigma(\mathbf{w}_\ell^\top \mathbf{h}_\ell + b_\ell)
\end{equation}
where $\mathbf{h}_\ell \in \mathbb{R}^d$ is the hidden representation at layer $\ell$.
We evaluate with AUROC (discrimination) and ECE (calibration).

\paragraph{Zero-Shot Cross-Lingual Transfer.}
The English-trained probe is applied directly to target-language representations $\mathbf{h}_\ell^{(\text{tgt})}$ without adaptation.

\paragraph{Transfer Conditions.}
We evaluate five transfer methods: (1)~zero-shot, (2)~mean-shift (center target representations to English mean), (3)~linear projection with $n$ parallel pairs, (4)~few-shot with $k$ target-language labels, (5)~CaCLR.

\paragraph{Mean-Shift Invariance.}
Under StandardScaler preprocessing and logistic regression, mean-shift adds a constant offset to all logits.
Since AUROC depends only on rank ordering, $\Delta\text{AUROC} \equiv 0$ mathematically.

\subsection{CaCLR: Calibration-Aware Cross-Lingual Regularization}
\label{sec:caclr_method}

CaCLR jointly optimizes English probe quality and cross-lingual consistency:
\begin{equation}
\mathcal{L}_{\text{CaCLR}} = \mathcal{L}_{\text{BCE}} + \alpha \cdot \mathcal{L}_{\text{SB-ECE}} + \beta \cdot \mathcal{L}_{\text{MMD}}
\end{equation}
where $\mathcal{L}_{\text{SB-ECE}}$ is smooth-binned ECE (differentiable calibration regularizer) and $\mathcal{L}_{\text{MMD}}$ is maximum mean discrepancy between English and target-language probe logits computed from parallel inputs:
\begin{equation}
\mathcal{L}_{\text{MMD}} = \text{MMD}^2_k\big(f(\mathbf{h}_\ell^{(\text{en})}), f(\mathbf{h}_\ell^{(\text{tgt})})\big)
\end{equation}
using Gaussian RBF kernel with median heuristic bandwidth.
Note that MMD operates on 1-D probe logits $f(\mathbf{h}) = \mathbf{w}^\top \mathbf{h} + b$, a design limitation: it can only match marginal distributions of scalar outputs, not the high-dimensional geometric structure of representations.

\subsection{Difficulty-Capability Decomposition}
\label{sec:decomp_method}

We decompose the correctness signal into difficulty-shared (language-agnostic) and capability-specific (language-dependent) components.
An IRT difficulty probe $\mathbf{w}_{\text{diff}}$ is trained to predict item difficulty from English representations.
The difficulty direction is projected out:
\begin{equation}
\tilde{\mathbf{h}}_\ell = \mathbf{h}_\ell - \frac{\mathbf{h}_\ell \cdot \mathbf{w}_{\text{diff}}}{\|\mathbf{w}_{\text{diff}}\|^2} \mathbf{w}_{\text{diff}}
\end{equation}
A correctness probe retrained on residuals $\tilde{\mathbf{h}}_\ell$ should transfer better if difficulty is the shared component and capability is the language-specific component driving degradation.

\subsection{Linear Projection}
\label{sec:linproj_method}

Given $n$ parallel pairs $\{(\mathbf{h}_\ell^{(\text{en})}_i, \mathbf{h}_\ell^{(\text{tgt})}_i)\}$, we learn a linear mapping $\mathbf{W}$ minimizing:
\begin{equation}
\min_{\mathbf{W}} \sum_i \|\mathbf{W} \mathbf{h}_\ell^{(\text{tgt})}_i - \mathbf{h}_\ell^{(\text{en})}_i\|^2
\end{equation}
This directly estimates the rotation correction between source and target representation spaces.

% ============================================================
\section{Experimental Setup}
\label{sec:setup}

\paragraph{Models.}
Qwen3-235B-A22B (94 layers, MoE, INT4 quantization), Llama~4 Scout (48 layers, MoE, W4A16), and Gemma~4 31B Dense (60 layers, FP16).
Quantization validation (Section~\ref{sec:quant}) confirms $\leq$1.1pp AUROC difference vs.\ FP16.

\paragraph{Languages.}
20 languages spanning 3 resource tiers: \textbf{High} (zh, es, fr, de, ja, pt, nl, ru), \textbf{Mid} (ar, hi, ko, vi, tr, id, pl, uk), \textbf{Low} (sw, bn, te).
Camera-ready extension covers 52 languages.

\paragraph{Benchmark.}
GlobalMMLU test split, partitioned into probe\_train (70\%) and eval (30\%) via stratified split.
English BCE labels and parallel activations for MMD come from probe\_train only.
Validation on Belebele and MMLU-ProX.

\paragraph{Probed Layers.}
Qwen3: 24 layers (every 4th); Llama~4 Scout: 12 layers (every 4th); Gemma~4: 15 layers (every 4th).

\paragraph{CaCLR Implementation.}
Gaussian RBF kernel with median heuristic bandwidth; Adam optimizer ($\text{lr} = 0.001$, cosine decay); 50 epochs with early stopping (patience 10 on val BCE); initialized from standard LR solution; batch size 256.
Average parallel set size per language: 9,800 items (min: 4,200 for Telugu, max: 14,042 for English).

\paragraph{Statistical Testing.}
All comparisons use bootstrap 95\% CIs (1,000 resamples), paired permutation tests, and Cohen's $d$.

% ============================================================
\section{Results}

\subsection{Pareto Frontier Discovery}
\label{sec:pareto}

\paragraph{Qwen3-235B: Uniform Strong Tradeoff.}
Table~\ref{tab:qwen_transfer} summarizes the depth-transferability tradeoff on Qwen3.
English AUROC peaks at layer 72 ($\approx$77\% depth) with AUROC = 0.80.
However, the best \emph{transfer} layer is L48 ($\approx$51\% depth), where average zero-shot AUROC across 19 languages is 0.669 vs.\ 0.647 at L72---a 2.2pp gap.
All 20 languages exhibit negative Spearman $\rho$ between depth and transfer AUROC ($\rho_{\text{median}} = -0.48$, range $[-0.58, -0.42]$), confirming a universal tradeoff on this architecture.
Low-resource languages show stronger negative correlation (sw: $-0.58$, bn: $-0.56$, te: $-0.55$) due to weaker cross-lingual representations.

\begin{table}[t]
\centering
\small
\begin{adjustbox}{max width=\columnwidth}
\begin{tabular}{lcccc}
\toprule
& \multicolumn{2}{c}{\textbf{ZS AUROC}} & \multirow{2}{*}{\textbf{$\rho$}} \\
\cmidrule(lr){2-3}
\textbf{Language} & \textbf{L48} & \textbf{L72} & \\
\midrule
zh & 0.73 & 0.71 & $-0.45$ \\
es & 0.72 & 0.70 & $-0.42$ \\
fr & 0.71 & 0.70 & $-0.43$ \\
de & 0.70 & 0.69 & $-0.44$ \\
ja & 0.69 & 0.67 & $-0.47$ \\
pt & 0.70 & 0.69 & $-0.43$ \\
nl & 0.69 & 0.68 & $-0.44$ \\
ru & 0.69 & 0.67 & $-0.46$ \\
\midrule
ar & 0.67 & 0.65 & $-0.50$ \\
hi & 0.66 & 0.64 & $-0.52$ \\
ko & 0.68 & 0.66 & $-0.48$ \\
vi & 0.66 & 0.64 & $-0.51$ \\
tr & 0.65 & 0.63 & $-0.53$ \\
id & 0.67 & 0.65 & $-0.49$ \\
pl & 0.68 & 0.66 & $-0.46$ \\
uk & 0.67 & 0.65 & $-0.47$ \\
\midrule
sw & 0.58 & 0.55 & $-0.58$ \\
bn & 0.57 & 0.54 & $-0.56$ \\
te & 0.59 & 0.56 & $-0.55$ \\
\midrule
\textbf{Avg} & \textbf{0.669} & \textbf{0.647} & $-0.48$ \\
\bottomrule
\end{tabular}
\end{adjustbox}
\caption{Qwen3-235B zero-shot transfer AUROC at best transfer layer (L48, 51\% depth) vs.\ EN-optimal layer (L72, 77\% depth), with per-language Spearman $\rho$ (depth vs.\ transfer AUROC). All $\rho$ values are negative. Languages grouped by resource tier (High/Mid/Low).}
\label{tab:qwen_transfer}
\end{table}

\paragraph{Llama 4 Scout: Language-Specific Tradeoff.}
Table~\ref{tab:arch_contrast} contrasts per-language $\rho$ across architectures.
Llama~4 Scout shows a qualitatively different pattern: global $\rho = -0.12$, with 4 languages exhibiting \emph{positive} $\rho$ (ar: $+0.15$, ko: $+0.12$, hi: $+0.08$, id: $+0.02$), 8 near zero, and only 8 significantly negative.
This means deeper layers improve transfer for some languages while degrading it for others---there is no single ``best transfer layer'' on this architecture.
Per-language best layers vary from L20 (sw, bn) to L40 (ar, ko), whereas on Qwen3 all languages converge near L48.

\begin{table}[t]
\centering
\small
\begin{adjustbox}{max width=\columnwidth}
\begin{tabular}{lrr}
\toprule
\textbf{Language} & \textbf{Qwen3 $\rho$} & \textbf{Llama4 $\rho$} \\
\midrule
zh & $-0.45$ & $-0.35$ \\
es & $-0.42$ & $-0.30$ \\
fr & $-0.43$ & $-0.28$ \\
de & $-0.44$ & $-0.32$ \\
ja & $-0.47$ & $-0.25$ \\
pt & $-0.43$ & $-0.27$ \\
nl & $-0.44$ & $-0.22$ \\
ru & $-0.46$ & $-0.18$ \\
\midrule
ar & $-0.50$ & $\mathbf{+0.15}$ \\
hi & $-0.52$ & $\mathbf{+0.08}$ \\
ko & $-0.48$ & $\mathbf{+0.12}$ \\
vi & $-0.51$ & $-0.05$ \\
tr & $-0.53$ & $-0.08$ \\
id & $-0.49$ & $\mathbf{+0.02}$ \\
pl & $-0.46$ & $-0.10$ \\
uk & $-0.47$ & $-0.05$ \\
\midrule
sw & $-0.58$ & $-0.42$ \\
bn & $-0.56$ & $-0.38$ \\
te & $-0.55$ & $-0.30$ \\
\midrule
\textbf{Median} & $-0.48$ & $-0.10$ \\
\textbf{Range} & [$-0.58$, $-0.42$] & [$-0.42$, $+0.15$] \\
\bottomrule
\end{tabular}
\end{adjustbox}
\caption{Architecture contrast: per-language Spearman $\rho$ (depth vs.\ zero-shot transfer AUROC). Qwen3 shows uniformly negative $\rho$; Llama~4 Scout shows language-specific behavior with 4 positive values (\textbf{bold}).}
\label{tab:arch_contrast}
\end{table}

\paragraph{Gemma 4 31B Dense.}
The dense architecture exhibits the strongest and most uniform tradeoff ($\rho_{\text{median}} = -0.55$, range $[-0.62, -0.48]$, $\sigma = 0.04$), consistent with the absence of MoE routing heterogeneity.
EN best AUROC is 0.79 at layer 42 (70\% depth); best transfer layer is L30 (50\% depth) with median zero-shot AUROC 0.67.
The CKA cliff is steeper (magnitude 0.18) than both MoE models, providing the strongest architecture contrast.

\subsection{CaCLR: A Principled Negative Result}
\label{sec:caclr}

\paragraph{$\beta$ Sweep.}
Table~\ref{tab:beta_sweep} presents the MMD consistency loss $\beta$ sweep across three orders of magnitude.
No value of $\beta$ achieves $p < 0.05$ or $d > 0.2$.
At low $\beta$ (0.001--0.01), MMD gradients are too weak to influence probe weights.
At moderate $\beta$ (0.03--0.1), MMD pushes toward marginal distribution alignment of 1-D logits without addressing correctness direction rotation.
At high $\beta$ (0.3--1.0), MMD dominates the loss, degrading discriminative power ($\Delta$AUROC $= -0.010$ at $\beta = 1.0$).

\begin{table}[t]
\centering
\small
\begin{tabular}{lrr}
\toprule
$\beta$ & $\Delta$AUROC & $\Delta$ECE \\
\midrule
\rowcolor{lightgray}
0 (ref.) & 0.000 & 0.000 \\
0.001 & +0.002 & $-0.001$ \\
0.003 & +0.001 & $-0.001$ \\
0.01 & +0.001 & 0.000 \\
0.03 & $-0.001$ & +0.001 \\
0.1 & $-0.002$ & +0.002 \\
0.3 & $-0.005$ & +0.003 \\
1.0 & $-0.010$ & +0.005 \\
\midrule
\multicolumn{3}{l}{\small Paired perm.\ $p = 0.38$, Cohen's $d = 0.07$} \\
\bottomrule
\end{tabular}
\caption{$\beta$ sweep (MMD consistency weight). $\Delta$AUROC and $\Delta$ECE are relative to $\beta = 0$ (i.e., $\alpha$-only baseline with $\alpha = 0.5$). No $\beta$ value produces a significant effect.}
\label{tab:beta_sweep}
\end{table}

\paragraph{Component Ablation.}
Table~\ref{tab:ablation} decomposes CaCLR into its components.
The $\alpha$ (SB-ECE) term reduces ECE from 0.175 to 0.125 ($-$5.0pp) with minimal AUROC impact (+0.004pp).
The $\beta$ (MMD) term contributes $+$0.001pp AUROC and $-$0.001pp ECE---functionally zero.

\begin{table}[t]
\centering
\small
\begin{tabular}{lcc}
\toprule
\textbf{Configuration} & \textbf{ZS AUROC} & \textbf{ECE} \\
\midrule
BCE only & 0.647 & 0.175 \\
$+\alpha$ (SB-ECE) & 0.651 & 0.125 \\
$+\beta$ (MMD) & 0.649 & 0.173 \\
$+\alpha + \beta$ (Full) & 0.652 & 0.124 \\
\bottomrule
\end{tabular}
\caption{CaCLR component ablation (Qwen3, L48). The $\alpha$ term drives ECE improvement; $\beta$ contributes nothing.}
\label{tab:ablation}
\end{table}

\paragraph{Full Method Comparison.}
Table~\ref{tab:method_comparison} ranks all transfer methods.
Linear projection with 50 pairs (+3.5pp on Qwen3, +3.2pp on Llama~4) and few-shot with 100 labels (+4.6pp, +4.4pp) significantly outperform zero-shot transfer.
CaCLR achieves only +0.8pp on Qwen3 and +0.6pp on Llama~4, almost entirely from the $\alpha$ component.
Mean-shift produces $\Delta$AUROC $\equiv 0$ as predicted by the rank-invariance proof.

\begin{table}[t]
\centering
\small
\begin{adjustbox}{max width=\columnwidth}
\begin{tabular}{lcccc}
\toprule
& \multicolumn{2}{c}{\textbf{Qwen3 (L48)}} & \multicolumn{2}{c}{\textbf{Llama4 (L28)}} \\
\cmidrule(lr){2-3} \cmidrule(lr){4-5}
\textbf{Method} & AUROC & $\Delta$ & AUROC & $\Delta$ \\
\midrule
Zero-shot & 0.669 & --- & 0.638 & --- \\
Mean-shift & 0.669 & 0.000 & 0.638 & 0.000 \\
CaCLR ($\alpha$ only) & 0.673 & +0.004 & 0.642 & +0.004 \\
CaCLR (full) & 0.677 & +0.008 & 0.644 & +0.006 \\
\midrule
Linear proj.\ (50) & 0.704 & +0.035 & 0.670 & +0.032 \\
Few-shot (25) & 0.695 & +0.026 & 0.662 & +0.024 \\
Few-shot (100) & 0.715 & +0.046 & 0.682 & +0.044 \\
\bottomrule
\end{tabular}
\end{adjustbox}
\caption{Transfer method comparison. Linear projection and few-shot methods significantly outperform zero-shot and CaCLR. $\Delta$ is vs.\ zero-shot. Bootstrap 95\% CI for linear proj.\ $\Delta$: $[0.025, 0.045]$ (Qwen3).}
\label{tab:method_comparison}
\end{table}

\subsection{CKA Geometric Analysis}
\label{sec:cka}

CKA (linear kernel) between English and 19 target languages declines monotonically with depth: median CKA drops from 0.83 at shallow layers to 0.48 at deep layers.
At 70\% depth, resource-tier banding is clear (High: 0.58, Mid: 0.46, Low: 0.34).
Qwen3 exhibits a steeper CKA cliff at 72\% depth (magnitude 0.15) than Llama~4 at 65\% depth (0.08), consistent with the stronger Pareto tradeoff.

\paragraph{Partial Correlation Analysis.}
Table~\ref{tab:cka_partial} presents the key statistical test.
Pooled CKA--transfer correlation is $r = 0.52$ ($p < 10^{-6}$), but this is confounded by depth.
Partial correlation controlling for normalized depth drops to $r = 0.18$ ($p = 0.11$)---CKA adds marginal information beyond depth itself.
We therefore frame CKA as \emph{suggestive correlational analysis}, not geometric explanation.

The between-language claim remains valid: at matched depth, languages with higher CKA show less transfer degradation ($r = 0.52$, $p = 0.02$, $n = 19$).

\begin{table}[t]
\centering
\small
\begin{tabular}{lccl}
\toprule
\textbf{Analysis} & $r$ & $p$ & \textbf{Interpretation} \\
\midrule
CKA--transfer (pooled) & 0.52 & $<10^{-6}$ & Strong but confounded \\
Depth--transfer & --- & --- & Primary driver \\
CKA $|$ depth (partial) & 0.18 & 0.11 & Suggestive only \\
\midrule
Between-lang (same depth) & 0.52 & 0.02 & Valid cross-lang claim \\
\bottomrule
\end{tabular}
\caption{CKA as transfer predictor. Partial correlation controlling depth reduces $r$ from 0.52 to 0.18, indicating CKA's predictive power is largely driven by depth. Between-language correlation at matched depth remains significant.}
\label{tab:cka_partial}
\end{table}

\paragraph{Subspace Angle Analysis.}
Correctness subspace angles between English and target languages increase from median 15\textdegree{} at shallow layers to 38\textdegree{} at deep layers, far below the random baseline of 72\textdegree.
This quantifies the rotation of correctness directions with depth: the subspace persists but rotates, consistent with the Pareto frontier---deeper layers encode stronger English correctness but rotate away from cross-lingual alignment.

\subsection{Difficulty-Capability Decomposition}
\label{sec:decomp}

Table~\ref{tab:decomp} summarizes the decomposition ablation.
Removing the difficulty direction before retraining correctness probes improves transfer in 65\% of language-layer combinations (paired permutation $p = 0.028$, Cohen's $d = 0.25$).
The effect is concentrated at deep layers (78\% improved, $d \approx 0.35$) and high-resource languages (72\% improved), with diminishing returns for low-resource languages (50\% improved) where probes are noisier.
A random-direction control confirms specificity: random removal yields $d = 0.05$ ($p = 0.42$).

\begin{table}[t]
\centering
\small
\begin{tabular}{lccc}
\toprule
\textbf{Condition} & \textbf{Med.\ $\Delta$AUC} & \textbf{\% Impr.} & $d$ \\
\midrule
Difficulty removed & +0.015 & 65\% & 0.25 \\
Random direction & +0.003 & 52\% & 0.05 \\
\midrule
\multicolumn{4}{l}{\small \textit{Difficulty removed, by depth:}} \\
\quad Shallow (20--40\%) & +0.005 & 48\% & --- \\
\quad Mid (40--60\%) & +0.015 & 65\% & --- \\
\quad Deep (60--90\%) & +0.022 & 78\% & --- \\
\midrule
\multicolumn{4}{l}{\small \textit{Difficulty removed, by resource tier:}} \\
\quad High & +0.020 & 72\% & --- \\
\quad Mid & +0.013 & 63\% & --- \\
\quad Low & +0.008 & 50\% & --- \\
\bottomrule
\end{tabular}
\caption{Decomposition ablation. Removing the difficulty direction improves transfer, concentrated at deep layers and high-resource languages. Random-direction control confirms specificity ($p = 0.42$). Effect size is small ($d = 0.25$), supporting a qualitative rather than quantitative claim.}
\label{tab:decomp}
\end{table}

The IRT difficulty probe achieves median $R^2 = 0.40$ (range $[0.22, 0.52]$), peaking at 65\% depth.
The difficulty--correctness cosine of 0.20 (range $[0.10, 0.32]$) indicates incomplete separation---the decomposition is informative but not clean.

\subsection{Cross-Lingual Orthogonality}
\label{sec:ortho}

Table~\ref{tab:ortho} extends the correctness-confidence orthogonality finding to 20 languages.
All cosine values remain below 0.12 (English: 0.032 on Qwen3, 0.036 on Llama~4), confirming that output-based calibration is structurally limited regardless of language.
A resource-tier gradient emerges: high-resource median 0.037, mid 0.054, low 0.090 (Spearman $\rho = -0.58$, $p = 0.008$ with log resource estimate).
Cross-architecture correlation is 0.94, indicating the pattern is model-agnostic.

\begin{table}[t]
\centering
\small
\begin{adjustbox}{max width=\columnwidth}
\begin{tabular}{lcc|lcc}
\toprule
\textbf{Lang} & \textbf{Qwen3} & \textbf{Llama4} & \textbf{Lang} & \textbf{Qwen3} & \textbf{Llama4} \\
\midrule
en & .032 & .036 & ar & .055 & .060 \\
zh & .038 & .042 & hi & .058 & .065 \\
es & .035 & .038 & ko & .050 & .055 \\
fr & .033 & .036 & vi & .060 & .068 \\
de & .036 & .040 & tr & .062 & .070 \\
ja & .040 & .045 & id & .048 & .052 \\
pt & .034 & .037 & pl & .040 & .045 \\
nl & .033 & .036 & uk & .045 & .050 \\
ru & .042 & .048 & sw & .095 & .110 \\
 &  &  & bn & .088 & .098 \\
 &  &  & te & .078 & .085 \\
\midrule
\multicolumn{3}{l}{\textbf{Tier medians:}} & \multicolumn{3}{l}{High: .037, Mid: .054, Low: .090} \\
\bottomrule
\end{tabular}
\end{adjustbox}
\caption{Correctness-confidence cosine similarity across 20 languages and 2 architectures. All values $< 0.12$. Resource-tier gradient: $\rho = -0.58$, $p = 0.008$.}
\label{tab:ortho}
\end{table}

\subsection{Linear Projection Sample Efficiency}
\label{sec:linproj}

Table~\ref{tab:linproj} presents the linear projection learning curve.
With only 10 parallel pairs, median $\Delta$AUROC is +1.2pp; at 50 pairs, it reaches +3.5pp (90\% of achievable improvement).
Diminishing returns set in around 60 pairs (knee), consistent with cross-lingual alignment literature.
Linear projection surpasses CaCLR at only 20 pairs.
Low-resource languages benefit most (+4.5pp at 50 pairs vs.\ +2.8pp for high-resource), as their larger geometric misalignment provides more room for correction.

\begin{table}[t]
\centering
\small
\begin{tabular}{rccl}
\toprule
\textbf{Pairs} & \textbf{Med.\ $\Delta$AUC} & \textbf{Std} & \textbf{Note} \\
\midrule
10 & +0.012 & 0.022 & Below CaCLR \\
20 & +0.023 & 0.016 & Matches CaCLR \\
50 & +0.035 & 0.011 & 90\% of max \\
100 & +0.040 & 0.008 & Near saturation \\
200 & +0.042 & 0.007 & Saturated \\
\midrule
\multicolumn{4}{l}{\small CaCLR surpassed at ${\sim}$20 pairs; knee at ${\sim}$60} \\
\bottomrule
\end{tabular}
\caption{Linear projection sample efficiency. $\sim$50 parallel pairs achieve 90\% of maximum improvement. CaCLR is surpassed at $\sim$20 pairs.}
\label{tab:linproj}
\end{table}

\subsection{Control Task Analysis}
\label{sec:control}

Table~\ref{tab:control} presents the Hewitt \& Liang control task analysis extended to 20 languages.
Approximately 27\% of language-layer combinations exhibit selectivity $< 0.05$, concentrated along two dimensions: shallow layers universally (85\% below threshold at 0--20\% depth) and low-resource languages at mid-depth (42\% below for sw/bn/te).

\begin{table}[t]
\centering
\small
\begin{tabular}{lccc}
\toprule
& \textbf{Shallow} & \textbf{Mid} & \textbf{Deep} \\
& (0--30\%) & (30--60\%) & (60--100\%) \\
\midrule
High resource & 80\% & 12\% & 5\% \\
Mid resource & 88\% & 22\% & 8\% \\
Low resource & 92\% & 38\% & 15\% \\
\midrule
\textbf{All} & 85\% & 18\% & 8\% \\
\bottomrule
\end{tabular}
\caption{Fraction of language-layer combinations with selectivity $< 0.05$ (capacity artifact risk). Total: ${\sim}$27\%. Shallow layers and low-resource languages are most affected.}
\label{tab:control}
\end{table}

This has two implications.
First, shallow-layer low selectivity is expected and benign---the paper does not claim correctness signal at shallow layers.
Second, low-resource mid-depth artifacts are a genuine concern: probes with selectivity $< 0.05$ may reflect capacity memorization rather than correctness signal.
We recommend selectivity $\geq 0.05$ as a minimum threshold for interpretable cross-lingual probing.
Our own statistical tests exclude combinations below this threshold where feasible.

\subsection{Probe Complexity Validation}
\label{sec:probe_complexity}

MLP probes (2-layer) achieve +2.2pp English AUROC over logistic regression (0.825 vs.\ 0.80 on Qwen3; 0.80 vs.\ 0.78 on Llama~4) but suffer a 5.4pp larger transfer gap (18.5\% vs.\ 13.1\% on Qwen3).
Logistic regression transfers better for 82\% of languages; exceptions are zh and es where MLP nonlinearity partially transfers.
This validates the focus on linear probes: the transferable component of correctness signal is geometric (linear subspace), and nonlinear components are language-specific.

\subsection{Additional Validation}
\label{sec:quant}

\paragraph{Quantization.}
Comparing W4A16 vs.\ FP16 on Llama~4 Scout (5 layers, 5\% subset) yields max $\Delta$AUROC of 0.009 with Pearson $r = 0.996$.
GPTQ-Int4 vs.\ FP16 on Qwen3 yields max $\Delta$AUROC of 0.011 ($r = 0.995$).
All differences are within bootstrap CIs.

\paragraph{Cross-Benchmark Generalization.}
On Belebele (reading comprehension), Pareto $\rho = -0.45$ and linear projection $\Delta = +3.0$pp.
On MMLU-ProX (translated MCQA), $\rho = -0.48$ and $\Delta = +3.4$pp.
Method ranking is consistent across all three benchmarks (rank correlation with GlobalMMLU: 0.80 for Belebele, 0.90 for MMLU-ProX).
CaCLR $\beta$ remains inert on all benchmarks.

\paragraph{Extended Language Coverage.}
Expanding to 52 languages, median zero-shot AUROC is 0.64 (range $[0.50, 0.73]$).
10 languages fall below 0.55 (am, my, zu, sn, rw, xh, ny, mg), reflecting minimal pretraining representation.
With linear projection, 67\% of languages (35/52) achieve viable calibration ($\text{AUROC} > 0.60$).

% ============================================================
\section{Discussion}

\paragraph{Why Does MMD Fail?}
MMD operates on 1-D probe logits and can only match marginal distributions (mean and variance of scalar outputs).
Cross-lingual correctness misalignment is \emph{geometric rotation}---the correctness direction points differently in different languages' activation spaces.
This rotation occurs in high-dimensional space \emph{before} the probe, where MMD on post-probe logits cannot reach.
Zhao et al.\ (\citeyear{zhao2019}) predict exactly this: marginal matching fails when $P(Y|X)$ differs across domains, as correctness label distributions are language-dependent.
Linear projection succeeds because it directly estimates the rotation matrix via parallel pairs.

\paragraph{Architecture Dependence.}
With $n = 3$ architectures differing in type (MoE vs.\ dense), scale, layer count, and training data, we cannot attribute tradeoff shape to specific architectural features.
The finding is observational: practitioners must empirically characterize the Pareto frontier for their target architecture rather than assuming a universal depth rule.

\paragraph{Limitations.}
(1)~MCQA only: clean correctness labeling requires closed-form evaluation; open-ended generation requires different signal types.
(2)~$n = 3$ architectures: insufficient for causal attribution.
(3)~Decomposition effect size ($d = 0.25$) supports a qualitative claim only.
(4)~CKA is correlational after controlling for depth.
(5)~27\% of probing combinations are potential capacity artifacts, particularly for low-resource languages.

% ============================================================
\section{Conclusion}

We present the first systematic study of cross-lingual correctness probe transfer, revealing an architecture-dependent depth-transferability Pareto frontier.
The natural distributional-matching approach (CaCLR with MMD consistency) fails---a principled negative result explained by the geometric nature of cross-lingual misalignment.
The practical recommendation is simple: linear projection with ${\sim}50$ parallel pairs provides +3.5pp AUROC, surpassing CaCLR at 20 pairs and achieving 90\% of the maximum achievable improvement.
Cross-lingual orthogonality between correctness and confidence persists across all 20 tested languages (max cosine 0.11), establishing output-based calibration as structurally limited regardless of language.
We further contribute a methodological warning: approximately 27\% of cross-lingual probing results are capacity artifacts, and recommend selectivity $\geq 0.05$ as a minimum interpretability threshold.

% ============================================================
% References
\bibliographystyle{plainnat}

\begin{thebibliography}{30}

\bibitem[Angelopoulos and Bates,2023]{angelopoulos2023conformal}
Anastasios Angelopoulos and Stephen Bates. 2023.
\newblock Conformal prediction: A gentle introduction.
\newblock {\em Foundations and Trends in Machine Learning}, 16(4):494--591.

\bibitem[Artetxe et~al.,2018]{artetxe2018vecmap}
Mikel Artetxe, Gorka Labaka, and Eneko Agirre. 2018.
\newblock A robust self-learning method for fully unsupervised cross-lingual mappings of word embeddings.
\newblock In {\em Proceedings of ACL}, pages 789--798.

\bibitem[Ben-David et~al.,2010]{bendavid2010}
Shai Ben-David, John Blitzer, Koby Crammer, Alex Kulesza, Fernando Pereira, and Jennifer~Wortman Vaughan. 2010.
\newblock A theory of learning from different domains.
\newblock {\em Machine Learning}, 79(1):151--175.

\bibitem[Conneau et~al.,2018]{conneau2018muse}
Alexis Conneau, Guillaume Lample, Marc'Aurelio Ranzato, Ludovic Denoyer, and Herv{\'e} J{\'e}gou. 2018.
\newblock Word translation without parallel data.
\newblock In {\em Proceedings of ICLR}.

\bibitem[Hewitt and Liang,2019]{hewitt2019control}
John Hewitt and Percy Liang. 2019.
\newblock Designing and interpreting probes with control tasks.
\newblock In {\em Proceedings of EMNLP}, pages 2733--2743.

\bibitem[Kornblith et~al.,2019]{kornblith2019cka}
Simon Kornblith, Mohammad Norouzi, Honglak Lee, and Geoffrey Hinton. 2019.
\newblock Similarity of neural network representations revisited.
\newblock In {\em Proceedings of ICML}, pages 3519--3529.

\bibitem[Miao et~al., 2026a]{miao2026coral}
Qingyi Miao, Mingxuan Ju, and Lyle Ungar. 2026a.
\newblock {CORAL}: Cross-task correctness probes for LLM calibration.
\newblock In {\em Proceedings of NAACL}.

\bibitem[Miao and Ungar,2026]{miao2026steering}
Qingyi Miao and Lyle Ungar. 2026.
\newblock Steering LLM confidence via orthogonal correctness directions.
\newblock In {\em Proceedings of ACL}.

\bibitem[Pires et~al.,2019]{pires2019multilingual}
Telmo Pires, Eva Schlinger, and Dan Garrette. 2019.
\newblock How multilingual is multilingual {BERT}?
\newblock In {\em Proceedings of ACL}, pages 4996--5001.

\bibitem[Raghu et~al.,2017]{raghu2017svcca}
Maithra Raghu, Justin Gilmer, Jason Yosinski, and Jascha Sohl-Dickstein. 2017.
\newblock {SVCCA}: Singular vector canonical correlation analysis for deep learning dynamics and interpretability.
\newblock In {\em Proceedings of NeurIPS}, pages 6076--6085.

\bibitem[Voita and Titov,2020]{voita2020mdl}
Elena Voita and Ivan Titov. 2020.
\newblock Information-theoretic probing with minimum description length.
\newblock In {\em Proceedings of EMNLP}, pages 183--196.

\bibitem[Wu and Dredze,2019]{wu2019beto}
Shijie Wu and Mark Dredze. 2019.
\newblock Beto, bentz, becas: The surprising cross-lingual effectiveness of {BERT}.
\newblock In {\em Proceedings of EMNLP}, pages 833--844.

\bibitem[Xiao et~al.,2025]{xiao2025}
Shuyue Xiao, Zhiwei Zhang, and Linyi Yang. 2025.
\newblock Cross-lingual knowledge boundary probing in large language models.
\newblock In {\em Proceedings of ACL}, pages 1204--1218.

\bibitem[Zhao et~al.,2019]{zhao2019}
Han Zhao, Remi Tachet des Combes, Kun Zhang, and Geoffrey Gordon. 2019.
\newblock On learning invariant representations for domain adaptation.
\newblock In {\em Proceedings of ICML}, pages 7523--7532.

\bibitem[Zhou et~al.,2024]{zhou2024lace}
Kaitlyn Zhou, Dan Jurafsky, and Tatsunori Hashimoto. 2024.
\newblock {LACE}: Language-aware cross-lingual calibration.
\newblock In {\em Proceedings of EMNLP}, pages 2891--2906.

\end{thebibliography}

% ============================================================
\appendix
\section{Supplementary Results}

\subsection{Probe Complexity (Table S1)}

Table~\ref{tab:probe_complexity_supp} compares logistic regression and MLP probes.
MLP achieves higher English AUROC but larger transfer gap at all layers, with the gap widening at deeper layers (6.5pp at deep vs.\ 2.0pp at shallow).

\begin{table}[h]
\centering
\small
\begin{tabular}{lcccc}
\toprule
& \multicolumn{2}{c}{\textbf{EN AUROC}} & \multicolumn{2}{c}{\textbf{Transfer Gap}} \\
\cmidrule(lr){2-3} \cmidrule(lr){4-5}
\textbf{Model} & LogReg & MLP & LogReg & MLP \\
\midrule
Qwen3 & 0.800 & 0.825 & 13.1\% & 18.5\% \\
Llama4 & 0.780 & 0.800 & 14.2\% & 19.5\% \\
\midrule
\multicolumn{5}{l}{\small LogReg better transfer: 82\% of langs} \\
\multicolumn{5}{l}{\small LogReg ECE: 0.135; MLP ECE: 0.148} \\
\bottomrule
\end{tabular}
\caption{Probe complexity comparison. MLP gains English AUROC at the cost of larger transfer gap.}
\label{tab:probe_complexity_supp}
\end{table}

\subsection{Cross-Benchmark Generalization (Table S2)}

\begin{table}[h]
\centering
\small
\begin{tabular}{lccccc}
\toprule
\textbf{Benchmark} & \textbf{EN} & \textbf{ZS} & $\rho$ & \textbf{LP $\Delta$} & $\beta$ $\Delta$ \\
\midrule
GlobalMMLU & 0.80 & 0.669 & $-0.50$ & +.035 & +.002 \\
Belebele & 0.77 & 0.650 & $-0.45$ & +.030 & +.002 \\
MMLU-ProX & 0.81 & 0.670 & $-0.48$ & +.034 & +.003 \\
\midrule
\multicolumn{6}{l}{\small Rank corr.\ with GlobalMMLU: Belebele 0.80, MMLU-ProX 0.90} \\
\bottomrule
\end{tabular}
\caption{Cross-benchmark generalization (Qwen3). EN = EN best AUROC; ZS = zero-shot median; LP = linear projection 50 pairs. CaCLR $\beta$ inert across all benchmarks.}
\label{tab:benchmark_supp}
\end{table}

\subsection{Quantization Validation (Table S3)}

\begin{table}[h]
\centering
\small
\begin{tabular}{lcccc}
\toprule
& \multicolumn{2}{c}{\textbf{Llama4 (W4A16)}} & \multicolumn{2}{c}{\textbf{Qwen3 (INT4)}} \\
\cmidrule(lr){2-3} \cmidrule(lr){4-5}
\textbf{Layer} & $\Delta$AUC & & $\Delta$AUC & \\
\midrule
L8 / L16 & $-0.004$ & & $-0.006$ & \\
L16 / L32 & $-0.004$ & & $-0.005$ & \\
L24 / L48 & $-0.005$ & & $-0.004$ & \\
L32 / L64 & $-0.007$ & & $-0.005$ & \\
L40 / L80 & $-0.009$ & & $-0.011$ & \\
\midrule
\textbf{Mean} & $-0.004$ & $r{=}.996$ & $-0.005$ & $r{=}.995$ \\
\bottomrule
\end{tabular}
\caption{Quantization validation. All $\Delta$AUROC within bootstrap CIs. Quantization does not distort correctness signals.}
\label{tab:quant_supp}
\end{table}

\subsection{Control Task Details (Table S4)}

English selectivity increases with depth on both architectures: shallow average 0.03, mid 0.15--0.18, deep 0.24--0.27.
Low-resource languages show lower selectivity at matched depth (median 0.09 vs.\ 0.19 for high-resource), with the worst combinations being sw at 0--50\% depth (median selectivity 0.02) and te at 0--40\% depth (0.02).
All control task AUROCs remain near chance (median 0.50, range $[0.48, 0.53]$), confirming negligible probe capacity for memorizing random labels.

\subsection{Three-Architecture Summary (Table CR1)}

\begin{table}[h]
\centering
\small
\begin{adjustbox}{max width=\columnwidth}
\begin{tabular}{lccccccc}
\toprule
\textbf{Model} & \textbf{Type} & \textbf{EN} & \textbf{ZS} & $\rho_{\text{med}}$ & $\rho_{\sigma}$ & \textbf{CKA cliff} & \textbf{LP $\Delta$} \\
\midrule
Qwen3-235B & MoE & .80 & .669 & $-.48$ & .05 & .15 & +.035 \\
Llama4 Scout & MoE & .78 & .638 & $-.10$ & .12 & .08 & +.032 \\
Gemma4 31B & Dense & .79 & .670 & $-.55$ & .04 & .18 & +.033 \\
\bottomrule
\end{tabular}
\end{adjustbox}
\caption{Three-architecture comparison. Dense model shows strongest, most uniform Pareto ($\rho_\sigma = 0.04$). MoE models vary: Qwen3 moderate uniform, Llama4 language-specific.}
\label{tab:three_arch}
\end{table}

\subsection{Extended 52-Language Coverage (Table CR2)}

Of 52 languages tested, 32 achieve zero-shot AUROC $> 0.60$ and 25 exceed 0.65.
With linear projection (+3.2pp median), 35/52 (67\%) reach viable calibration.
8 new languages fall below 0.55 (am, my, zu, sn, rw, xh, ny, mg), reflecting minimal pretraining representation.
The resource-tier gradient replicates at scale ($\rho_{\text{median}} = -0.47$ across 52 languages).

\end{document}
